%%%%%%%%%%%%%%%%%%%%%%%%%%%%%%%%%%%%%%%%%%%%%%%%%%%%%%%%%%%%%
\section{Introduction}\label{sec:intro}
%%%%%%%%%%%%%%%%%%%%%%%%%%%%%%%%%%%%%%%%%%%%%%%%%%%%%%%%%%%%%

This paper presents a quasi-monolithic cell-centred finite volume method for fluid--solid interaction (FSI), implemented in the OpenFOAM-based toolbox solids4foam~\citep{Cardiff2025} and solved using a Jacobian-free Newton--Krylov (JFNK) procedure.
The remainder of this section places the method in the context of existing FSI approaches, identifies the gap addressed, and summarises the specific contributions.

FSI methods can be classified along three axes: the modelling viewpoint (Eulerian, Lagrangian, or mixed), the degree of interface coupling (weak vs.\ strong), and the solver structure (partitioned vs.\ monolithic).

Eulerian approaches treat both solid and fluid domains as ``fluid-like'' (e.g.\ volume-of-fluid~\citep{jain2019conservative}, level-set~\citep{dunne2006eulerian, cottet2008eulerian}), easily allowing arbitrarily large solid deformations, but conservative interface descriptions and the inclusion of advanced solid constitutive laws remain challenging.
In contrast, Lagrangian (particle-type) approaches treat both domains as ``solid-like'', e.g.\ smooth-particle hydrodynamics, the material point method, and the particle finite element method (PFEM)~\citep{aubry2005particle, aubry2006fractional, cremonesi2020state}.
These methods are promising for large-deformation problems but typically come with a prohibitive computational cost.
Mixed Eulerian--Lagrangian approaches treat the solid (and interface) as Lagrangian and the fluid using an arbitrary Lagrangian--Eulerian (ALE) description.
ALE is effective for small to moderate interface deformations but eventually fails for very large deformations~\citep{wall2006large}, motivating methods that allow the solid to ``float over'' an Eulerian fluid grid, such as immersed boundary/domain~\citep{peskin1977numerical, peskin2002immersed, iaccarino2003immersed}, Chimera/overset~\citep{ge2003three, miller2014overset, li2021fluid, cane2021patient}, and cut-cell methods~\citep{pasquariello2016cut, tao2020new, xie2020cartesian}.
Of these, overset methods offer large-deformation freedom while retaining orthogonal interface cells suitable for resolving fluid boundary layers, albeit with the downside of greater implementation complexity.

In weakly coupled (loosely coupled) approaches, the fluid and solid are solved once per time step without iteration over the interface conditions~\citep{farhat2000two}.
This is computationally inexpensive but unstable for strongly coupled systems with significant added-mass effects~\citep{Causin2005, Forster2006}, such as incompressible flows interacting with light, flexible structures.
Strongly coupled approaches strictly enforce the interface kinematic and dynamic conditions at each time step; the present work focuses on the strongly coupled case.

Strong coupling can be enforced through either a partitioned or a monolithic solver structure.
In partitioned (staggered, two-system) approaches, separate solvers --- and potentially separate discretisations --- are used for the fluid and solid~\citep{Tukovic2018}, and outer iterations are performed over the interface at each time step.
These are accelerated using Aitken relaxation~\citep{Degroote2009, bungartz2016precice, irons1969version, kuttler2008fixed}, quasi-Newton schemes~\citep{santiago2020hpc, delaisse2022surrogate}, or Robin-type coupling~\citep{Tukovic2019}.
The principal limitations are robustness and efficiency: for strongly coupled, highly nonlinear problems, outer-iteration convergence can be slow or fail entirely, particularly under strong added-mass loading.
In contrast, monolithic approaches treat the fluid and solid as a single system, solved concurrently within the same linearised system~\citep{Degroote2009}.
This avoids outer iterations and can handle tightly coupled nonlinear problems robustly, at the cost of reduced modularity and a poorly conditioned linear system that is more difficult to solve.

Monolithic FSI methods have been developed predominantly in the finite element (FE) community.
Early influential examples include \citet{Heil2004}, who first demonstrated that block-triangular approximations of the coupled Jacobian provide effective preconditioners for monolithic ALE FSI with GMRES, and the segregated/monolithic comparison of \citet{Heil2008}.
The Munich line of \citet{Mayr2015, Schott2019a} extended this with temporally consistent and cut-FEM-based formulations; in particular, \citet{Mayr2015} provides a clear FE counterpart to block-preconditioned monolithic Newton--Krylov FSI through Newton with GMRES and per-field ILU(0) preconditioning.
The benchmark of \citet{Turek2006} (used as a test case in Section~\ref{sec:test_cases}) and the broader analyses of \citet{Richter2012, Richter2017} provided rigorous verification and goal-oriented frameworks.
\citet{Liu2014} introduced a quasi-monolithic FE formulation specifically targeting strong added-mass effects, in which the fluid mesh motion is treated outside the coupled Newton vector --- the approach most closely related to the present work.
A comprehensive treatment is available in the textbook of \citet{JaimanJoshi2022}.
% and the recent open-source FEniCS-based turtleFSI solver~\citep{Bergersen2020, Gjertsen2017}.
An early matrix-free FE precedent is the PFEM-based monolithic Lagrangian method of \citet{Ryzhakov2010}, in which global pressure condensation enables a matrix-free diagonally-preconditioned conjugate gradient inner solve.
Most recently, fully coupled Newton--Krylov methods with exact Jacobians and multigrid/Schwarz preconditioning have been demonstrated in the FE setting~\citep{Aulisa2018}, alongside ongoing development of field-split and block preconditioners for added-mass-dominated FSI~\citep{Badia2008, Muddle2012, Deparis2016, Jodlbauer2019, Calandrini2020}.
A common feature across this body of work is that the underlying spatial discretisation is finite element.

In contrast, body-fitted cell-centred finite volume FSI methods --- particularly within the OpenFOAM ecosystem --- have largely followed partitioned coupling strategies.
The standard reference here is \citet{Tukovic2018}, whose partitioned framework underpins the solids4foam toolbox \citep{Cardiff2025} and a substantial body of subsequent applied work; a recent review~\citep{Cardiff2021} documents the broader trajectory of finite volume solid mechanics that this work builds on.
Cross-ecosystem coupling frameworks such as preCICE~\citep{bungartz2016precice} extend this paradigm but remain partitioned in nature.
Vertex-centred FV variants such as the hybrid scheme of \citet{Malan2013} achieve strong coupling through partitioned dual-time iteration but, by the authors' own description, remain block-Jacobi at the FSI level rather than monolithic.
A smaller line of cell-centred FV work instead adopts unified or one-system formulations, beginning with \citet{Greenshields1999b} and \citet{Greenshields2005} and later continued by \citet{Giannopapa2006, Giannopapa2008} and \citet{Tandis2023}.
These unified FV-FSI procedures solve common pressure--velocity-type equations across fluid and solid regions using segregated SIMPLE-type pressure--velocity splitting, rather than a fully coupled Newton or Newton--Krylov solve.
Separate Eulerian FV one-continuum approaches have also emerged, notably the quadtree-based monolithic method of \citet{Bergmann2022} and recent VOF-based work~\citep{Prajapati2025}, in which fluid and structure phases are solved over a common Cartesian computational domain without distinct fluid and solid meshes.

The closest prior work in terms of solver architecture is the parallel monolithic FV/FE scheme of \citet{Eken2016}, which couples a side-centred (staggered) unstructured finite volume fluid discretisation with a Galerkin finite element solid using inexact Newton with FGMRES, an upper-triangular right block preconditioner replacing the zero pressure block with a scaled discrete Laplacian, and restricted additive Schwarz subdomain solves via PETSc.
The present method differs in three important respects: the fluid discretisation is cell-centred collocated rather than staggered, with both fluid and solid in the same cell-centred FV framework; the fluid mesh motion is treated outside the coupled Newton vector (a quasi-monolithic strategy) rather than included as additional unknowns; and the Newton--Krylov operator is applied in a Jacobian-free manner, with the approximate block preconditioner serving solely as the preconditioning matrix rather than as an assembled Jacobian.

To the authors' knowledge, no prior work has presented a body-fitted cell-centred finite volume quasi-monolithic FSI formulation solved with a Jacobian-free Newton--Krylov procedure and an approximate fluid--solid block preconditioner; the closest precedent is the staggered-FV/FE monolithic Newton--Krylov method of \citet{Eken2016}.
The present work fills this gap, building on the authors' recent JFNK methods for cell-centred FV solid mechanics~\citep{Cardiff2026, Batistic2026}, which are here extended to coupled FSI.

The specific contributions are:
\begin{enumerate}
    \item A cell-centred finite volume formulation for incompressible Newtonian fluid--elastic solid interaction, in which the fluid momentum (in ALE form), fluid pressure, and solid momentum (in total Lagrangian form) equations are coupled within a single nonlinear residual, with strong enforcement of the interface kinematic and dynamic conditions. The fluid mesh motion is treated outside the coupled Newton vector, following the quasi-monolithic strategy of \citet{Liu2014}.
    \item A Jacobian-free Newton--Krylov solution strategy supported by an approximate block preconditioner. The preconditioner is assembled from compact FV linearisations of the fluid, solid, and interface-coupling terms, and supports both monolithic and nested field-split usage --- outer fluid--solid split, with optional inner velocity--pressure Schur-complement approximation via the Rhie--Chow pressure-stabilisation block.
    \item An open-source implementation in the solids4foam toolbox~\citep{Cardiff2025} interfaced with PETSc~\citep{PETSc}, demonstrating that monolithic FSI is achievable within a segregated finite volume framework without rewriting the underlying OpenFOAM discretisation.
    The source code and benchmark cases are released openly, supporting independent verification, reproduction, and extension by the community.
    \item A benchmark study covering established 2D and 3D FSI test cases, including verification of spatial and temporal accuracy, comparison against partitioned solutions, and assessment of solver robustness under strong added-mass conditions.
\end{enumerate}

Section~\ref{sec:math_model} presents the governing equations for the coupled fluid--solid problem.
Section~\ref{sec:numerical_methods} develops the cell-centred FV discretisation, the JFNK solver, and the approximate block preconditioner; the fluid mesh-motion procedure and implementation in solids4foam/PETSc are also described.
Section~\ref{sec:test_cases} presents numerical results across eight benchmark configurations.
Section~\ref{sec:conclusion} concludes.
\hl{Before submission, verify all references have correct details}